\defmodule {LATIO}

This module contains several tools for input/output,
for performing generator searches, for managing and storing lists
of generators, and the like.
It contains the actual implementations of the
high-level ``seek'' and ``lat'' programs.
The procedures here use most of the other 
lower-level and intermediate-level modules.
All the {\tt Read...} procedures read the data in the format
explained in the documentation of the {\tt seekl} program.

The searches for lists of (good or bad) generators are 
performed by categories, as explained in
the documentation of the {\tt seekl} program.
The set of dimensions for which the lattice structure is examined
is partitioned into a certain number of intervals (one or more),
also called {\em categories\/}.
We give bounds on the accepted values of the figure of merit
for each category and, for any given category $c$,
we consider only those generators that satisfy the bounds 
for all the categories less or equal to $c$.
Then we retain a list of generators with the best or worst figures of
merit for each category.


\bigskip\hrule
\code

DEFINITION MODULE LATIO;

FROM SUPINT   IMPORT SuperInteger;
FROM SUPFACT  IMPORT Factors;
FROM CONFIG   IMPORT BDim, OutputType, MaxJ, MaxCat, Indexj, MaxNbGen;
FROM MULT     IMPORT MScal, MVect;
FROM NORM     IMPORT CriterionType; 
FROM REDBAS   IMPORT VerifyType; 
IMPORT BASISLR;


TYPE
   GenerType   = (MRG, MWC);
\endcode
 \tabb Type of generator that can be considered. \endtabb
\code

   CondType    = (NoCond, AppFact, MaxBits);
\endcode
 \tabb Type of condition on the multipliers. \endtabb
\code

   VerifType   = (VerifyAll2, VerifyRetained2, 
                  VerifyAll3, VerifyRetained3, NoVerify);
\endcode
 \tabb Type of verification that can be performed on the results
    (i.e., type of computational method). \endtabb
\code

   Poly        = POINTER TO InfoPoly;
   InfoPoly    = RECORD
      a        : MVect;
      NextPoly : Poly
      END;
\endcode
 \tabb  Elements to store a linked list of polynomials.  \endtabb
\code

   Resultats   = POINTER TO InfoRes;
   InfoRes     = RECORD
      RD2, RS2, RQ2 : BASISLR.BVect;
      DimD2, DimQ2  : BDim
      END;
\endcode
 \tabb  Type of record used as a field of the record type {\tt InfoGen},
    to store detailed results for a given generator. 
    It contains the {\em squared\/} values of the $d_t$, $S_t$, and $q_t$,
    and the maximal dimensions {\tt DimD2} and {\tt DimQ2} 
    for which $d_t$ and $q_t$, respectively, have been computed.
 \endtabb
\code

   Generateur = POINTER TO InfoGen;
   InfoGen    = RECORD
      Maj      : ARRAY Indexj OF MVect; 
      a        : MVect;
      DimMerit : BDim; 
      Merit    : LONGREAL;
      Res      : Resultats;
      NextGen  : Generateur
      END;
\endcode
 \tabb  Type of record used to store a linked list of generators.
   Used to store the generators that are kept in each category.
   The coefficients $a_{j,i}$ of the components are in {\tt Maj}
   whereas those of the combined generator are in {\tt a}.
   {\tt DimMerit} is the dimension where the minimum is attained for
   the figure of merit (for the given category) and {\tt Merit} gives
   the value of this figure of merit.
   {\tt Res} is a pointer to the more detailed results and
   {\tt NextGen} is a pointer to the next generator in the list.
 \endtabb
\code

   ToZone = POINTER TO Zone; 
   Zone   = RECORD
      No         : INTEGER; 
      Binf, Bsup : SuperInteger;
      p          : LONGREAL;
      PetiteDiff : BOOLEAN;
      BsupMsH    : SuperInteger;
      Mq, MBsup  : MScal;
      NextZone   : ToZone;
      END;
\endcode
 \tabb  Type of record used to store the research zones.
   {\tt No} is the number of the zone, which is in $\{0,1,2,3\}$.
   {\tt Binf} and {\tt Bsup} are the lower and upper bounds
   defining the zone.  
   {\tt P} is the fraction of the acceptable
   values of $a$ lying in this zone (this is used for the random searches).
   {\tt PetiteDiff} is {\tt TRUE} iff {\tt Bsup - Binf $\le$ H} (or {\tt Hk}),
   and when it is {\tt TRUE}, {\tt BsupMsH} is equal to 
   {\tt Bsup ${}+1-{}$H} (or {\tt Hk}). This is also used for random searches.
   The variables {\tt Mq} and {\tt MBsup} are used in {\tt ExamAllZones}.
  \endtabb
\code

   ToRegion = POINTER TO Region;
   Region   = RECORD
      No         : INTEGER;
      Binf, Bsup : MScal;
      Mq         : MScal
      END;
\endcode
 \tabb  Type of record used to store the regions that are to be searched.
   {\tt No} is the zone number where this region is (in $\{0,1,2,3\}$).
   {\tt Binf} and {\tt Bsup} are the lower and upper bounds
   defining the region.  
   The variable {\tt Mq} is  used in {\tt ExamRegion}.
  \endtabb
\code

PROCEDURE TimeOver () : BOOLEAN;
\endcode
 \tab 
   Returns {\tt TRUE} if and only if the cpu time limit has been reached.
 \endtab
\code

PROCEDURE CreateGenerateur (VAR G : Generateur);
\endcode
 \tab 
  Creates a record for a generator {\tt G} and initializes its fields.
 \endtab
\code

PROCEDURE CreateZone (VAR Z : ToZone);
\endcode
 \tab 
  Creates a record for a zone {\tt Z} and initializes its  fields.
 \endtab
\code

PROCEDURE CreateRegion (VAR R : ToRegion);
\endcode
 \tab 
  Creates a record for a region {\tt R} and initializes its  fields.
 \endtab
\code

PROCEDURE LireFact (VAR L : Factors;   NomFch : ARRAY OF CHAR);
\endcode
 \tab
  Reads a list of factors in the file {\tt NomFich} and places them
  in the list {\tt L}.  Assumes that the factors in the file are
  written according to the format used by {\tt Ecrirefact}.
 \endtab
\code

PROCEDURE EcrireFact (L : Factors;   NomFch : ARRAY OF CHAR);
\endcode
 \tab
  Writes the factors of the list {\tt L} to the file {\tt NomFich}.
  It this file exists, it is rewritten.
  The factors can be read back later by {\tt LireFact}.
 \endtab
\code

PROCEDURE Readm (j : Indexj);
\endcode
 \tab
  Reads the modulo $m_j$ for the component $j$.
  Initializes some constants that depend on it.
 \endtab
\code

PROCEDURE ReadOnea (aa : SuperInteger);
\endcode
 \tab
  Reads one multiplier $a_{j,i}$ or one bound 
  $b_{j,i}$ or $c_{j,i}$ in the format described in {\tt seekl}.
 \endtab
\code

PROCEDURE ReadReadGen;
\endcode
 \tab
  Reads a boolean indicating if the generators to be tested should be
  taken from a `{\tt .gen}' file, and if true, \
  also reads the name of the file, without the {\tt .gen} extension.
 \endtab
\code

PROCEDURE ReadTypeGen (j : Indexj);
\endcode
 \tab
  Reads the type of generator to be considered, MRG or MWC,
  for component $j$.
 \endtab
\code

PROCEDURE ReadOrder (j : Indexj);
\endcode
 \tab
  Reads the order $k_j$ of the recurrence  for component $j$.
 \endtab
\code

PROCEDURE ReadPerMax (j : Indexj);
\endcode
 \tab
  Reads a boolean indicating if component $j$ should have 
  maximal period.  The muliplier $m_j$ must have been read in
  the `3 numbers' format by {\tt readmj}, otherwise this boolean
  is reset to {\tt FALSE}.
 \endtab
\code

PROCEDURE ReadImplemCond (j : Indexj);
\endcode
 \tab
  Reads the condition that we impose on the coefficients $a_{j,i}$
  for the component $j$.  
 \endtab
\code

PROCEDURE ReadFactmr (j : Indexj);
\endcode
 \tab
  Reads what to do about the factorizations of $m_j-1$ and 
  $r_j = (m_j^{k_j}-1)/(m_j-1)$,  for the component $j$.
 \endtab
\code

PROCEDURE Reada (j : Indexj);
\endcode
 \tab
  Reads the $k_j$ multipliers $a_{j,1},\dots,a_{j,k_j}$ for component $j$.
 \endtab
\code

PROCEDURE Readbc (j : Indexj);
\endcode
 \tab
  Reads the $2k_j$ bounds $b_{j,1},c_{j,1}\dots,b_{j,k_j},c_{j,k_j}$ 
  on the multipliers for component $j$.
 \endtab
\code

PROCEDURE ReadSearchMethod (j : Indexj);
\endcode
 \tab
  Reads the method of search, exhaustive or random, for the polynomials
  for component $j$.  If the method is {\tt Random}, the it reads also
  the number of regions and their sizes.
 \endtab
\code

PROCEDURE ReadMRGComponent (j : Indexj);
\endcode
 \tab
  Reads all the parameters related to the component $j$,
  by calling the appropriate `{\tt Read...}' procedures above.
  Also creates the required data structures for this component.
 \endtab
\code

PROCEDURE ReadMinMaxDim;
\endcode
 \tab
  Reads the minimal and maximal numbers of dimensions in which the
  test is to be applied, the number of categories,
  and the boundaries of these categories.
  The format is slightly different for the `seek' and `lat'
  programs.
 \endtab
\code

PROCEDURE ReadMinMaxMerit;
\endcode
 \tab
  Reads the minimum and maximum acceptable values for the figure 
  of merit for each category of dimensions.
 \endtab
\code

PROCEDURE ReadNbGen;
\endcode
 \tab
  Reads the number of generators that we want to retain in each
  category.
 \endtab
\code

PROCEDURE ReadCriterion;
\endcode
 \tab
  Reads the selection criterion (figure of merit).
 \endtab
\code

PROCEDURE ReadLatInfo;
\endcode
 \tab
  Reads what lattice information we want to be printed in the results.
 \endtab
\code

PROCEDURE ReadLatType;
\endcode
 \tab
  Reads which type of lattice we want to analyze:
  That generated by the recurrent states only, or that generated by
  a single orbit, or that generated by the longest cycle of a 
  generator with prime power modulus, or the full lattice generated
  by the union of all the orbits.
  In the case of {\tt Orbit}, this procedure also reads the initial
  state and initializes some additional data structures.
 \endtab
\code

PROCEDURE ReadLaIndices;
\endcode
 \tab
  Reads the indices (lacunary or successive) that are to be considered.
 \endtab
\code

PROCEDURE ReadVerifyBB;
\endcode
 \tab
  Reads the verification level that we want to use.
 \endtab
\code

PROCEDURE ReadMaxNodesBB;
\endcode
 \tab
  Reads the maximum number of nodes that will be allowed in
  the branch-and-bound tree.  When that number is exceeded,
  the branch-and-bound procedure will stop.
 \endtab
\code

PROCEDURE ReadTimeLimit;
\endcode
 \tab
  Reads the cpu time limit that is allowed for a search.
  When that limit is exceeded, the search procedure stops
  and reports what it has found so far.
 \endtab
\code

PROCEDURE ReadSeeds;
\endcode
 \tab
  Reads the seeds that are used for the random number generator
  of SENTIERS.  This generator is used for the random searches.
 \endtab
\code

PROCEDURE ReadResultForm;
\endcode
 \tab
  Reads the form of the output (on the terminal, in a {\tt .gen} file,
  in a {\tt .res} file, in a {\tt .tex} file).
 \endtab
\code

PROCEDURE Computem;
\endcode
 \tab  Computes the order and the modulus $m$ for the 
   combined generator, given those of the components.
   Computes also the period lengths of the components and of
   the combination.
 \endtab
\code

PROCEDURE Computea;
\endcode
 \tab 
   Computes the vector of multipliers $a$ for the combined generators,
   given the parameters of the components.
   Assumes that the moduli of the components are all prime.
   Initializes the vectors {\tt Maj} and {\tt a}. 
 \endtab
\code

PROCEDURE SaisieDonnees;
\endcode
 \tab
  Calls the different low-level procedures to read the data
  from the data file, in the format appropriate for the programs
  {\tt latl} and {\tt seekl}.
  Also creates some of the data structures (e.g., {\tt SuperInteger}'s)
  that are required.
 \endtab
\code

PROCEDURE ReadGenFile;
\endcode
 \tab
  Reads a set of generators from the file `{\tt xxx.gen}',
  assuming that the data has been read from `{\tt xxx.dat}'. 
 \endtab
\code

PROCEDURE PerMaxPowPrime (j : Indexj) : BOOLEAN;
\endcode
 \tab
   Assumes that $m_j$ is a power of a prime $p$, that $k_j = 1$, 
   and that the recurrence is homogeneous.   
   Returns {\tt TRUE} iff the maximal period conditions are 
   satisfied for that case for component $j$.
 \endtab
\code

PROCEDURE Reducemj (VAR a : MVect);
\endcode
 \tab
  Reduces $m_j$ and reinitializes the module {\tt MRG} accordingly
  in the case where the modulus is a power of a prime and where
  {\tt LatType = PrimePower}.  Called by {\tt testGen}.
 \endtab
\code

PROCEDURE InitZones (j, i : INTEGER);
\endcode
 \tab
  Initializes the research zones for the multiplier $a_i$ of the
  component $j$.
  In the case of a random search, this procedure also creates the
  region for $(j,i)$.    Called by {\tt SetupGenSeek}.
 \endtab
\code

PROCEDURE ChoisirBornes (j : Indexj);
\endcode
  \tab  This procedure is called by {\tt SeekGen} in the case of a
    random search.  It chooses a region at random and 
    initializes its boundaries.  This region will then be searched
    completely.  
  \endtab
\code

PROCEDURE TestGen
   ( VAR a    : MVect;
     T        : CriterionType;
     VerifyBB : VerifyType;
     Out      : OutputType
   );
\endcode
 \tab
   Tests the (combined) generator with multiplier (vector) {\tt a}.
 \endtab
\code

PROCEDURE ConserverGen (C : INTEGER;  tMerit2 : BDim);
\endcode
 \tab
   Inserts the most recently tested generator into the list
   for the categoty {\tt Cat}.
   The value {\tt tMerit2} indicates the dimension where the
   worst-case is attained for the figure of merit.
   This procedure is called by {\tt VerifyCategories}.
 \endtab
\code

PROCEDURE VerifyCategories (VAR Me2 : BASISLR.BVect);
\endcode
 \tab
  Verifies if the most recently tested generator can be included
  in any of the lists corresponding to the different categories.
  If so, includes it in the appropriate lists. 
  The vector {\tt Me2} is either {\tt Q2} or {\tt S2}, depending
  on the figure of merit that is used.
 \endtab
\code

PROCEDURE ExamAllZones (j, i : INTEGER);
\endcode
 \tab
  This procedure is used in the ``exhaustive search'' case. 
  When $J > 1$, it collects primitive polynomials for the given $j$.
  When $J = 1$, it tests all generators in the region.
  The initial call should be with $i$ equal to the degree of the
  polynomials for component $j$.
  Then, the procedure calls itself recursively for smaller $i$.
 \endtab
\code

PROCEDURE ExamRegion (j, i : INTEGER);
\endcode
 \tab
  Similar to {\tt ExamAllZones}, but used for the ``Random search'' case.
 \endtab
\code

PROCEDURE ExamCombPoly (jj : Indexj);
\endcode
 \tab
  Used only when $J > 1$.  Examines all combinations of the primitive
  polynomials found by {\tt ExamRegion} or {\tt ExamAllZones}.
 \endtab 
\code

PROCEDURE ExamCurrenta;
\endcode
 \tab
  Tests the generator that corresponds to the current value of the
  multiplier (vector) {\tt a}, and adds it to the list(s) if it is competitive.
  Called only by {\tt TestGenFromFile}.
 \endtab
\code

PROCEDURE TestGenFromFile;
\endcode
 \tab
  Examines all the generators from a {\tt .gen} file,
  which is assumed to be the current input file.
 \endtab
\code

PROCEDURE CalculNbreGen (j : Indexj);
\endcode
 \tab
  Computes the total number of simple generators (i.e., polynomials)
  examined, for the component {\tt j}, and puts the result in
  {\tt TotGen [j]}.
 \endtab
\code

PROCEDURE CalculerResteInfo;
\endcode
 \tab 
  After the search has been completed, computes all the remaining required
  information on the retained generators.
  Computes the $d_t$ or the $q_t$ when required and not already done.
  Redo the computations with error bounds if requested.
 \endtab
\code

PROCEDURE PrintErrBounds (T : INTEGER);
\endcode
 \tab
  Computes and prints the bounds on the Euclidean distances between
  the points produced by a L'Ecuyer-style combined MRG and the 
  corresponding points produced by the associated MRG,
  in dimensions 2 to {\tt T}.  This is used only when $J > 1$.
 \endtab
\code

PROCEDURE EcrisGen (G : Generateur;    Cat : INTEGER);
\endcode
 \tab
  Writes the parameters and the results relative to a retained 
  generator {\tt G}, in category {\tt Cat}.
  Called by {\tt WriteSeekRes}.
 \endtab
\code

PROCEDURE WriteLatHead;
\endcode
 \tab
  Writes a heading for the results of a lattice test.
  Called by {\tt LatGen}.
 \endtab
\code

PROCEDURE WriteLatHeadTex;
\endcode
 \tab
  Writes a heading for the results of a lattice test in \LaTeX\ form.
  Called by {\tt LatGen}.
 \endtab
\code

PROCEDURE WriteDataComponent (j : Indexj);
\endcode
 \tab
  Write the statistical data relative to component {\tt j}.
 \endtab
\code

PROCEDURE WriteStatComponent (j : Indexj);
\endcode
 \tab
   Writes the statistics on the searches for primitive 
   (or examined) polynomials for the component {\tt j}.
   Called by {\tt WriteSeekRes}.
 \endtab
\code

PROCEDURE WriteTabBest (l : LONGINT;  G : Generateur);
\endcode
 \tab
  Prints a one-line summary of the figures of merit in each interval,
  for best generator retained in category (i.e., interval) number {\tt l}.
 \endtab
\code

PROCEDURE WriteSeekHead;
\endcode
 \tab
  Writes, in the current output file, the heading and the data 
  for a search by the {\tt SeekGen} procedure.
  Writes also the factorizations in the files if this was requested.
 \endtab
\code 

PROCEDURE WriteSeekRes;
\endcode
 \tab
   Writes the results of a search in the current output file,
   then closes it.
 \endtab
\code

PROCEDURE WriteSeekGen;
\endcode
 \tab
   Writes in a {\tt .gen} file the generators found in a search.
   If the data was read in {\tt xxx.dat}, then this procedure
   will open the file {\tt xxx.gen} and write in it.
 \endtab
\code

PROCEDURE SetupSeekGen;
\endcode
 \tab
  Called by {\tt SetUp}.
  Performs initializations for the searches of generators.
  Creates the ordered lists to store the `best' {\tt NbGen [C]}
  generators in each category {\tt C}.
  Initializes the bounds on the figures of merit in each dimension.
  Initializes the polynomial lists (to empty) and the search zones.
 \endtab
\code

PROCEDURE Setup;
\endcode
 \tab
  Obtains and opens data file, reads data and performs initializations.
  Initializes {\tt REDBAS} and {\tt MRG}.
  This procedure is called by {\tt SeekGen} and {\tt LatGen}.
 \endtab
\code

PROCEDURE LatGen;
\endcode
 \tab  Applies the required tests 
   (for ``lat'' programs) and prints results.
 \endtab
\code

PROCEDURE SeekGen;
\endcode
 \tab  Performs a seek for ``good'' generators 
   (implements ``seek'' programs).
 \endtab
\code

PROCEDURE CalcLaIndicesStreams (s, w : LONGINT);
\endcode
 \tab Calculates lacunary indices by groups of $s$ spaced by $2^w$
 \endtab 
\code

PROCEDURE LatGenStreams;
\endcode
 \tab  Applies the test for different values of $w$ where a big generator
  is divided into smaller ones, which are separated by $2^w$. 
 \endtab 
\code 

END LATIO.
\endcode








